% ============================================================
%  Poul Martin Møller, "Qvindelighed"
%
%  Written 1822--1823 (per the editor's note); published
%  posthumously in: Efterladte Skrifter af Poul M. Møller,
%  bd. 5 (Kjøbenhavn: C.A. Reitzel, 1855--56), pp. 5--13.
%  Transcribed from that edition's Fraktur setting.
%
%  TRANSCRIPTION (Danish)
%  Conventions: original orthography kept (aa-spellings,
%  capitalised nouns, æ/ø). Fraktur long-s rendered as s.
%  Letterspaced emphasis in the original -> \emph{}.
%  The editor's illegibility mark "[uløseligt]" is reproduced
%  as printed. Printed page numbers -> \opage{N}.
% ============================================================

\documentclass[12pt]{article}
\usepackage[utf8]{inputenc}
\usepackage[T1]{fontenc}
\usepackage{libertinus}
\usepackage{libertinust1math}
\usepackage[danish]{babel}
\usepackage[protrusion=true,expansion=true]{microtype}
\usepackage[margin=1.4in]{geometry}
\usepackage{setspace}
\usepackage{fancyhdr}
\usepackage{marginnote}
\usepackage[colorlinks=true,linkcolor=black,urlcolor=black,
  pdftitle={Qvindelighed},
  pdfauthor={Poul Martin Møller}]{hyperref}
\setlength{\headheight}{15pt}
\pagestyle{fancy}
\fancyhf{}
\fancyhead[LE,RO]{\thepage}
\fancyhead[RE]{\textit{Qvindelighed}}
\fancyhead[LO]{\textit{Poul Martin Møller}}
\setlength{\parindent}{1.5em}
\setlength{\parskip}{0pt}
\onehalfspacing

% Original printed page numbers as marginal notes
\newcommand{\opage}[1]{\marginnote{\footnotesize\textit{[#1]}}}
% Centred rule marking the author's section breaks
\newcommand{\qvbreak}{\par\medskip\begin{center}\rule{0.32\textwidth}{0.4pt}\end{center}\medskip}

\title{\textbf{Qvindelighed}}
\author{Poul Martin Møller}
\date{Skrevet 1822--1823; udgivet posthumt i \textit{Efterladte Skrifter},
  bd.~5,\\ Kjøbenhavn: C.A.\ Reitzel, 1855--56, pp.~5--13}

\begin{document}
\maketitle

\opage{5}For blandt Qvinder at ansees som et Mønster paa Fortræffelighed
behøver man ikke saa meget at udføre noget Dygtigt i Verden, som at yttre høi
Begeistring for Dyden og føre skjønne Grundsætninger paa Læben. Ogsaa i denne
Henseende ere Damer mere spiritualistiske; de gjøre ikke nogen Fordring paa, at
Dyden --- i det Mindste ikke deres egen --- skal yttre sig i de grove
materialistiske Handlinger. En ædlere Natur bringer dem til at holde sig
bestandig i Tankens høiere Region, de elske de sædelige Begreber aldeles for
deres egen Skyld, uden at kræve dem afprægede i det grove jordiske Stof. Dyden
taber egentlig Intet derved; thi hvad den forliser i Intension, det vinder den
igjen i Extension. De dydige Handlinger, en Dødelig i sit korte Jordliv kan
udføre, beløbe sig til et høist ringe Tal. Derimod kan en Danneqvinde blot i een
smuk Aftenstund have flere sværmende Forsætter og yndige Følelser, end Aristides
brugte sin hele Levetid. Og dog er Dyden ligefuldt tilstede, men blot som et
finere Væsen, opløst og \ldots\ [uløseligt]. Min salig Cousine var et mærkeligt
Exempel herpaa. Hendes gamle Moster, som tilbragte sine sidste Levedage hos
hende, havde faaet det Huusraad mod sin Gigt, at hun hver Aften skulde gnides paa
Ryggen med en Flonels Klud. Hun havde den Særhed, at Ingen uden min Cousine kunde
forrette denne Tjeneste ret efter hendes Hoved. Men uheldigviis pleiede denne
just til samme Tid at være fordybet i Fortællinger om Ridderfruer og Ædlinger,
der opoffrede deres jordiske Lykke for Guds Riges Skyld. Denne henrivende Læsning
\opage{6}kunde hun sjelden overtale sig til at ombytte med saa plump en
Haandgjerning, og Moster blev derfor almindeligviis forsømt. Engang sagde jeg til
hende: Gode Cousine! Du skulde dog forbarme dig lidt meer over det gamle Skrog.
Ifølge hendes Alder vil hun ikke længe kunne falde dig til Besvær, og naar du har
mistet hende, vil du sikkert fortryde din Mangel paa Opmærksomhed for hende. ---
Min Ven! svarede hun, du begriber ikke, hvor slige smaalige Sysler ere mig
modbydelige. Den Slags Selvovervindelse, som saa liden en Kjærlighedsgjerning
forudsætter, falder det mig svært at nedlade mig til. Jeg kan trods Nogen
fornegte mig selv for min lidende Søster og Broder, men kun i det Store; hvad
Enhver kan gjøre, det krymper jeg mig ved at give mig ifærd med. Mit Hjerte vil
altid blive miskjendt, fordi jeg er født i en ussel Tid, hvor jeg ikke har
hjemme, i en Tid, da Qvinden blot skal beskjæftige sig med Smaating. Døe for min
gode Moster kan jeg med Glæde, ja udgyde mit Blod for hende paa Retterstedet. Det
haaber jeg, du troer om mig; men at spilde min Tid for at føie hendes
Egensindighed, det holder jeg mig for god til.

Min Cousine følte sig virkelig i Stand til at døe for sin Moster, ja næsten for
Enhver. Men dermed vil jeg slet ikke have sagt, at denne høie Phantasie virkelig
ved given Leilighed her i denne Verden var blevet iværksat. Naar virkelig en
skiden Rakkerknegt, hvis Støvler lugtede af Tran, med uvasket Haand og uklippede
Negle havde grebet hende ved det blonde Hovedhaar, da havde hun sikkert i Tide
fattet en bedre Beslutning. Det var egentlig den rene Forestilling: at døe for
sin Næste, det fine Begreb, blottet for alle raae Tilsætninger, der begeistrede
hendes Sind. Hun var ei blevet sit Forsæt tro, hvis hun havde seet Dødsscenen
forestillet efter Naturen i sort Kobberstik, langt mindre, hvis Billedet havde
været illumineret.

Hun skyede saa meget som muligt al Berøring med den \opage{7}prosaiske
Virkelighed; thi Erfaringen gjorde daglig en ny Erobring fra det Tankerige, hvori
hun udførte sine ideelle Heltegjerninger. I en eenlig Stue, paa sin bløde Sopha,
naar hun ret uforstyrret overlod sig til Drømmen om, hvad hun under andre
Omstændigheder vilde have gjort, da drev hun sine Fordringer til Mennesket saa
høit, at hun med Grund foragtede Leonidas, Winkelried og Jeanne d'Arc; thi hun
vidste, at hun i deres Sted vilde have gjort Meer, end de gjorde. Blot Mucius
Scævola havde hun en hemmelig Respect for; hun følte, han var hendes Overmand.
Grunden dertil var, at hun engang havde brændt sig paa et Krusejern, og havde
derfor helligen lovet sig selv, aldrig frivilligen at komme Flammen for nær.
Derved var hendes Phantasieverden blevet indskrænket til tre Elementer. Men var
hun mindre stormodig, fordi hendes Helbred ikke tillod, at hun stod i
Vexelvirkning med det Synlige? Dersom sædeligt Værd skulde bedømmes efter den
Dygtighed, hvormed man forstaaer at udsætte sig for Uleilighed af ydre Ting, da
maatte jo Fiskerkonen, som sidder fra Morgen til Aften i Regn og Slud ved sin
Kurv, uden at faae Tandpine, faae særdeles høi Plads paa Moralitetens Rangliste.
Naar Enhver virker i den Sphære, hvori han føler sig hjemme, da kan man ikke
forlange Meer af ham. Een føler sig i Stand til at tumle med vaade Torsk og
Villinger, en Anden med Tanker og Forestillinger, eftersom Enhvers Constitution
tillader det. Men fristes man ikke ved slige Betragtninger til at ønske, at
Skaberen havde, til Brug for finere Væsner, beredt en finere Skueplads af
destillerede Elementer; at han havde besørget en beqvemmere Lommeudgave af Verden
til Brug for nervesvage Personer?

Jeg berørte før min Cousines Fordringer til sin Næstes Moralitet. Herover maa jeg
erklære mig noget bestemtere. Hun var i denne Henseende meget billig, og
forlangte ikke stort Meer af dem, end af sig selv. Hun krævede ei af Næsten, at
han \opage{8}skulde være streng i Opfyldelsen af andre Pligter end dem, han havde
mod hende. Forresten lod hun sig fuldkommen tilfredsstille, naar han blot havde
en smuk Livstheorie. Ædle Grundsætninger maatte han besidde, og have Talent til
at bringe dem for Dagen i et skjønt Sprog; da kaldtes han i hendes Terminologie:
en rar En. Om man endogsaa fortalte om en saadan Person, at han var en Dagdriver
og ikke gjorde for en Penning Gavn her i Verden, blev hun dog fast i sin
Hengivenhed for ham, og gav altid det undskyldende Svar: Ja, men han er dog
saadan en rar En; eller ogsaa: Ja, men han er dog saa sød.

En af vore fælleds Venner stod meget slet anskrevet hos min Cousine; thi han er
en stor Practicus, hvis Tid ikke tillader at give sig synderlig af med høie
Livsbetragtninger. Da han nødig vil udkaste nogen Plan, som han ikke seer sig i
Stand til at udføre, sætter han sig kraftigt imod alle høitstemte Fordringer til
de svage Dødelige. Følgen heraf var, at han gjerne temmelig kort affærdigede
Adelaide, naar hun søgte at drage ham ind i sin sværmende Tankekreds; ja med en
Slags Trodsighed yttrede han tidt en langt plattere Livsphilosophie, end den, han
selv følger. Denne Mand var hende naturligviis høist modbydelig. Jeg tog ham af
og til i Forsvar og bad hende lægge Mærke til, at han i Virkeligheden daglig viste
meer Selvopoffrelse, end de Fleste. Men her lagde hun igjen sin ideelle Natur for
Dagen, da hun satte Ordet og Begrebet høit over de plumpe Gjerninger. Han var og
blev hende reent uudstaaelig.

\qvbreak

I Almindelighed ivre Moralister strengt mod Qvinder, der give usædvanlig meget af
deres Legemsskjønhed til Priis for nysgjerrige Tilskueres Øine. Det klæder ikke
meget smukt; men der gives en anden Blottelse, som i Almindelighed roses, og som
dog i Grunden strider meget meer imod sand Qvindelighed. \opage{9}En aandrig
Qvinde har seet Sagen i det rette Lys, naar hun siger, at enhver sand Følelse
ledsages af en vis Blufærdighed. Denne Tro synes endeel af hendes Søstre slet
ikke at bekjende sig til, da de netop med Flid lægge deres smukke Følelser ret
tydeligen for Dagen og bruge Fremstillingen deraf som en Slags poetisk Prydelse
for deres Person. Men røber det Mangel paa Blufærdighed, ei tilbørligen at
bedække det ydre Menneske, da forudsætter det visseligen en endnu større
Tilsidesættelse af Blufærdighed, forsætligen at blotte det indre Menneske; at
gjøre Jagt paa Bifald ved en Udstilling af Hjertets ædle Beskaffenhed. --- En vis
Grad af Cultur gjør Følelsens ubevidste Udbrud høist sjeldne blandt Menneskene,
men den maa heller blive skjult i sit Lønkammer, end træde frem med selvbevidste
theatralske Phraser og Fagter. Naar en Qvinde med sit Vidende lægger sine indre
Bevægelser for Dagen, da tænker hun ifølge sin Natur nødvendigviis paa, om det
klæder hende smukt eller ikke. Deraf følger, at hendes Forhold ved slig Leilighed
bliver til et Slags fiint Koketterie. Hemmelig føler enhver Qvinde selv den sande
Betydning af en saadan Opførsel; thi naar hun ikke har megen Øvelse i at give
slige Forestillinger, da rødmer hun gjerne noget ved sin egen boglige Pathos. Men
en ret uskyldig Jomfru lægger fast aldrig sine Sindsbevægelser frem i Selskaber
af Flere; de skjules dybt i hendes Hjertes Løndom, eller udgydes for en Veninde
eller en Moder (ei for en Søster), og udvortes er hun omgivet af en tilsyneladende
Strenghed og Kulde, som kuns en Daare forvexler med Hjerteløshed.

\qvbreak

\begin{center}\large G.~F.~G.\end{center}

De maa ikke forundre Dem over, at jeg strax uden Indledning begynder paa Texten
til dette Brev. Min Hensigt er, \opage{10}nøiere at forklare et Par Ord, som jeg
dristigt, men løseligt, igaar Aftes i Samtalen med Dem og Deres Søster lod mig
forlyde med. Det er mig neppe muligt, at gjøre mig fuldkommen forstaaelig for
Dem; thi som en Følge af Tidens fragmentariske Udvikling bliver jo den Ene i vore
Dage næsten aldrig ret forstaaet af den Anden. Men jeg føler dog, at jeg muligen
kunde blive meget misforstaaet, naar jeg lod Sagen beroe ved vor mundtlige
Disput.

Først maa jeg da søge at godtgjøre, at jeg ingenlunde har villet komme Qvindens
Ære for nær ved at paastaae, hvad jeg paastaaer endnu, at det er \emph{uskjønt},
ja \emph{vederstyggeligt}, at et Fruentimmer er Digterinde af \emph{Profession}.
Dermed sætter jeg dem hverken over eller under Mænd, da jeg anseer det for høist
bagvendt, at ville anvende saadan Rangforordning paa vor Herres Have. Man kan i
det Høieste kun sammenligne Ting af samme Art; nærmere beseet, maaskee ei engang
det. Saameget bliver da vist, at man ikke kan bestemme, om f.~Ex. et Egetræ eller
Bøgetræ staaer øverst paa vor Herres Liste, om et Tordenskrald eller en Poppelpiil
er høiest \&c.\ \&c. Men det Naturværk bliver dog skjønnest, som kommer Naturens
Tanker nærmest, som renest fremstiller Skaberens Idee. Det er altsaa
uimodsigeligen smukkest, at en Qvinde ret er qvindelig. Naar man da paa sin Vei
igjennem Livet møder en Frøken***, da har man vist ikke Lov til at critisere den
bizarre Phantasie, som Naturen undertiden har overladt sig til i sære
Sammensætninger af stridige Elementer; men skjønt kalder man dog aldrig et saadant
Phænomen. Netop samme Virkning gjør det paa mig, at tænke mig et Fruentimmer som
lærd eller som --- Digterinde!! Jeg vil ikke appellere til Historien, skjøndt man
deraf lærer, at alle Poetricer have været aandelige Vanskabninger og Misfostre
(Sappho, Mad.\ Dacier, Mad.\ Stael, Fru B.\ \&c.). Men det støder høilig min
umiddelbare Tact for det Qvindeligskjønne, at see et Fruentimmer tygge \opage{11}%
paa sin Pen eller sidde med sin Foliant paa Skjødet. Derved bliver jeg netop
tilmode, som om jeg saae hende spille Trekort og trumfe dygtig i Bordet, bande høit
og dyrt, ryge sin Zigarro, gaae med Vandstøvler, samt med en rask mandhaftig
Spytten at beskrive store Buer i Luften. Der gives en vis activ Fremtræden i
Livet, hvortil jeg regner Forfatterskab, Digten, Præken, Slagsmaal, som i mine
Tanker klæde Fruentimmer ilde. Var De meget orthodox, kunde jeg beraabe mig paa
Apostelen Paulus, som siger, at Qvinden skal være taus i Kirken.

Betragte vi Poesiens Væsen nærmere, da bliver det endnu mere soleklart. Der gives
kun to Arter Digtning. Enten kan den historiske Erindring eller den nærværende Tid
afgive Stoffet. Det er vist den høieste Art af Kunst, der opstaaer som Efterklang
af en herlig forsvundet Tid, hvis Minde i smukke Sagn lever frisk paa Alles Læber.
Da maae naturligviis nogle særdeles Høitbegavede fremstaae, som renest opfatte
Billedet af Tidens sluttede Drama. Det sørger den Styrelse for, der ikke lader
noget Prægtigt forgaae uden Eftermæle. Denne Slags Kunst, hvori Historiens Fostre,
luttrede ved Tidens Skjærsild, atter træde frem med et fornyet Liv, er et reent
Naturproduct, ligesom de frodige Væxter, der myldre frem af de styrtede Ruiner.
Til denne Classe høre Kongerne for alle Digtere: Homer, Shakspeare og Ariost. Her
komme naturligviis Fruentimmerne slet ikke i Betragtning. Naar har en Qvinde
oplivet Oldtidens Erindring? Naar har hun blot \emph{forsøgt} paa at frembringe
noget i egentlig Forstand Høit eller Stort i Poesien?

Den anden Green af Poesien er den personlige. Det er den eneste virkelige Kunst,
der gives i vor Tid. (Thi de svage, halv hviskende Echoer fra Tider, der længst
have faaet deres kraftige Liigpræken, fortjene slet ikke at regnes med). Det
følger af sig selv, at den, hvis Liv har været rørt af stærke Lidenskaber, kommer,
naar Stormen har lagt sig, til Besindighed og kan, \opage{12}hvis han besidder
Fremstillingsgave, give os en interessant Fremstilling af, hvad han leed og
levede. Alle dygtige Digterværker i vor Tid, om de end komme nok saa forklædte
frem, have deres Rod i Forfatternes eget Levnet. Goethe erkjender selv, at Werther
og Faust ere Frugter af hans eget lidenskabelige Ungdomsliv, og kuns derved vinde
de den vidunderlig henrivende Sandhed. Ihvor stor en Aand Goethe er, har han
aldrig villet leve for at være Digter. Hans poetiske Værker ere under hans
mangehaande Studier, næsten uden Beslutning, uvilkaarligen komne til Verden.
Nuomstunder at leve for at være Digter, er det Samme som at reise for at gjøre
Reisebeskrivelse. Det lader altid ilde, at føle, for at skildre sine Følelser; men
dobbelt, naar Damer gjøre det. De kunne, som en Følge af Livets nuværende stramme
Former, ikke engang komme til at føre et saa broget Levnet, som den maa have
gjort, der skal have noget Mærkeligt at skildre. De, der have gjort det, ere ogsaa
Damer derefter. Den uskyldigste og almindeligste Damepoesie bestaaer i
Efterligning af, hvad Andre have skabt. Men slige Producter betyde dog i Kunstens
Rige slet ikke meer end de almindelige med Sortkridt tegnede Copier af Petrarch og
Laura, Achilles, Homer, Ossian \&c., som man seer hænge i et grønt Silkebaand paa
Kabinettets Vægge. Men hos de Qvinder, hos hvem Poesien virkelig er Noget, det vil
sige, hos hvem den er spundet ud af deres eget Bryst, er den langt mere
\emph{skyldig}. En saadan Digterinde maa ængstelig lure paa alle sine Følelser,
for strax at faae dem førte til Bogs. Ingen af hendes Glæder, ingen af hendes
Taarer gaaer tilspilde; Alt bliver net registreret i Riim og udstillet for
Publicum. Naar hun omfavner sin Moder, tænker hun paa den skjønne Attitude, hun
derved vinder for sit poetiske Pulterkammer; naar hun spadserer ud i Guds frie
Natur, har hun saa listelig en Smule Velinpapir i sin Sypose, for strax at have
Værktøiet ved Haanden, om hun skulde kunne overkomme saamegen Rørelse, som gaaer
\opage{13}paa et matthisonsk Digt. Bliver hendes Elsker troløs (som da vel hverken
Gud eller nogen Mand fortænker ham i), da glæder hun sig hemmelig over, at hun dog
derved profiterer en Snees Sonnetter eller nogle væsentlige Bestanddele til en
begrædelig Novelle. Kort sagt, hele Livet igjennem sidder hun for sit poetiske
Speil og silhouetterer sine smaae Følelser, saasnart de ere komne halv til Verden.

Dette er, gode Frue, min uforgribelige Mening om Digterinder, og i Følge den vil
det ikke forundre Dem, at jeg glæder mig over, at vi endda ikke have fleer af det
Slags Folk i vort Land. Nu vil jeg bede Dem ikke at blive vred, fordi jeg har
gjort Beslag paa Deres Tid med disse Linier. Jeg skulde ikke have vovet mig
dertil, naar jeg ikke af Deres Dom om Prækener var bragt til den Mening, at De
ikke var saa let at kjede. Alvorlig talt, vilde jeg gjerne udslette et af de
ubehagelige Indtryk, som jeg undertiden nødvendigviis maa gjøre paa Dem ved den
haarde, drøie, plumpe Maade, hvorpaa jeg saa tidt mod min Villie kommer til at
udtrykke mig hos Dem. Jeg beundrer virkelig den uendelige Blidhed og Godmodighed,
hvormed De finder Dem deri. Paa nogen Tid har jeg gandske mod min Villie faaet den
skarpe Fremtræden i Livet, som tidt bagefter endogsaa sætter mig selv i Forundring.

\begin{flushright}
Deres hengivne\\
P.~Møller.
\end{flushright}

\qvbreak

\end{document}
